% vim: ai sw=2 spell spelllang=cs

\begin{frame}{Typy souborů}

  \begin{itemize}
    \item Soubory s binárními daty
    \begin{itemize}
      \item Při zápisu a čtení nejsou data nijak měněna
    \end{itemize}

    \pause

    \item Soubory s textovými daty
    \begin{itemize}
      \item Binární data, ale interpretovaná
      \item Při čtení a zápisu můžou být některé bajty nahrazeny/změněny

      \pause

      \item Na Linuxu a jiných POSIX systémech ne
      \begin{itemize}
	\item Pozor na kompatibilitu s Windows!
	\item Na Windows se nahrazují konce řádků
      \end{itemize}
    \end{itemize}

    \pause

    \item Vlaječka při otevírání
  \end{itemize}

\end{frame}

\begin{frame}{Práce se soubory}{Pomocí stdlib}

  \begin{enumerate}
    \item Otevření souboru
    \begin{itemize}
      \item \code{fopen}
      \item Získáme držátko k souboru (\code{FILE *})
    \end{itemize}

    \pause

    \item Čteme/zapisujeme z/do souboru
    \begin{itemize}
      \item \code{fscanf}, \code{fprintf}, \code{fread}, \code{fwrite}, \dots
      \item Potřebují držátko k souboru
    \end{itemize}

    \pause

    \item Ukončíme práci se souborem (zavřeme soubor)
    \begin{itemize}
      \item \code{fclose}
    \end{itemize}

  \end{enumerate}

\end{frame}

\begin{frame}{Jak otevřít soubor -- mód otevření}

  \begin{itemize}
    \item Mód otevření volíme dle požadovaného chování (\doc{man fopen})

    \pause

    \item Čtení z existujícího souboru: \code{\"r\"}
    \item Vytvoření nového souboru a zápis do něj: \code{\"w\"}

    \pause

    \item Zápis na konec existujícího souboru: \code{\"a\"}

    \pause

    \item Čtení i zápis do nového souboru: \code{\"w+\"}

    \item Čtení i zápis do existujícího souboru
    \begin{itemize}
      \item Čtení i zápis kdekoli: \code{\"r+\"}
      \item Čtení kdekoli, zápis na konec: \code{\"a+\"}
    \end{itemize}

    \pause

    \item Práce v binárním režimu namísto textového: \code{\"b\"} na konec módu
    \begin{itemize}
      \item Např.\ \code{\"rb\"}
    \end{itemize}
  \end{itemize}

\end{frame}

\begin{frame}[fragile]{Otevření souboru}

  \begin{itemize}
    \item \code{FILE *fopen(const char *filename, const char *mode)}
    \begin{itemize}

      \item \code{filename} -- cesta k souboru
      \begin{itemize}
	\item Relativní: \file{test.txt}, \file{../test.txt}
	\item Absolutní: \file{/tmp/test.txt}
      \end{itemize}

      \pause

      \item \code{mode} -- mód otevření
      \begin{itemize}
	\item Dle předchozího slajdu
      \end{itemize}

      \pause

      \item Návratová hodnota -- držátko k souboru
      \begin{itemize}
	\item Při chybě: \code{NULL}
	\item Nemusíme znát definici \code{FILE} (interní záležitost OS)
      \end{itemize}
    \end{itemize}

    \pause

    \item Pozor na znak \code{\\} v řetězci cesty
    \begin{itemize}
      \item C pokládá \code{\\} za speciální znak, nutno použít escape sekvenci
	\code{\\\\}
      \item Např.\ soubor \file{test\\1.txt} $\rightarrow$
        \code{\"test\\\\1.txt\"}
    \end{itemize}

    \pause

    \item Zavření: \code{fclose(FILE *drzatko_z_open)}
  \end{itemize}

  \begin{block}{Příklad}
  \begin{lstlisting}
FILE *file = fopen("test.txt", "r");
fclose(file);
  \end{lstlisting}
  \end{block}

\end{frame}

\begin{frame}
  \frametitle{Úkol}

  \heading{Otevřete soubor}
  \begin{enumerate}
    \item Pracujte v adresáři \ifpcdirlinux \file{07} \else \file{11} \fi z
      repozitáře
    \item Otevřete pro čtení neexistující \file{text.txt}
    \begin{itemize}
      \item Ověřte návratovou hodnotu
    \end{itemize}

    \item Zavřete
    \item Vyzkoušejte
    \item Opravte, abyste otevírali správný \file{test.txt}
    \begin{itemize}
      \item Nutno nastavit správnou cestu aktuálního adresáře
      \item Vlevo \textit{Projects} $\rightarrow$ nahoře \textit{Run}
        $\rightarrow$ \textit{Working directory}
    \end{itemize}

    \item Znovu vyzkoušejte
  \end{enumerate}

\end{frame}

\begin{frame}[fragile]{\code{FILE} je platformně závislá}

  \begin{center}
  \begin{tabular}{|c|c|} \hline
  Linux & Windows \\ \hline
  \begin{lstlisting}
typedef struct _IO_FILE {
  int _flags;
  char *_IO_read_ptr;
  char *_IO_read_end;
  char *_IO_read_base;
  char *_IO_write_base;
  char *_IO_write_ptr;
  char *_IO_write_end;
  char *_IO_buf_base;
  char *_IO_buf_end;
...
} FILE;
  \end{lstlisting} &
  \begin{lstlisting}[showlines]
typedef struct _iobuf {
  char *_ptr;
  int _cnt;
  char *_base;
  int _flag;
  int _file;
  int _charbuf;
  int _bufsiz;
  char *_tmpfname;
} FILE;


  \end{lstlisting} \\ \hline
  \end{tabular}

  \bigskip
  \bigskip
  \pause

  \heading{Nespoléhejte na položky \code{FILE}!}
  \end{center}

\end{frame}

\begin{frame}{Čtení ze souboru}

  \begin{itemize}
    \item Čte se z aktuální pozice v souboru
    \begin{itemize}
      \item Po přečtení se pozice posune těsně za přečtená data
    \end{itemize}

    \pause

    \item Konec souboru
    \begin{itemize}
      \item Funkce vracejí \code{EOF} (\pozor, má typ \code{int})
      \item \code{feof} vrací nenulu
    \end{itemize}

    \pause

    \item Načtení jednoho znaku
    \begin{itemize}
      \item \code{int getc(FILE *stream)}
    \end{itemize}

    \pause

    \item Načtení jedné řádky (ukončené \code{\\n})
    \begin{itemize}
      \item \code{char *fgets(char *str, int num, FILE *stream)}
    \end{itemize}

    \pause

    \item Formátované čtení do proměnných
    \begin{itemize}
      \item \code{int fscanf(FILE *stream, const char *format, ...)}
    \end{itemize}

    \pause

    \item Blokové čtení na binární úrovni
    \begin{itemize}
      \item \code{size_t fread(void *ptr, size_t size, size_t count,
        FILE *stream)}
      \item Načte blok bajtů o zadané délce: \code{size * count}
    \end{itemize}
  \end{itemize}

\end{frame}

\begin{frame}[fragile]{Čtení ze souboru}{Příklad}

  \begin{block}{Čtení po znacích}
  \begin{lstlisting}
#include <err.h>
#include <stdio.h>

int main(void) {
  FILE *file;
  int value;

  file = fopen("test.txt", "r");
  if (!file)
    err(1, "fopen");

  while ((value = getc(file)) != EOF)
    putchar(value);

  fclose(file);
}
  \end{lstlisting}
  \end{block}

\end{frame}

\begin{frame}
  \frametitle{Úkol}

  \heading{Čtení souboru}
  \begin{enumerate}
    \item Otevřený soubor čtěte
    \item Vypište počet řádků v souboru
    \item Vyzkoušejte
  \end{enumerate}

\end{frame}

\begin{frame}{Zápis do souboru}

  \begin{itemize}
    \item Zapisuje se na aktuální pozici v souboru
    \begin{itemize}
      \item Po zápisu se pozice posune těsně za zapsaná data
    \end{itemize}

    \pause

    \item Zápis jednoho znaku
    \begin{itemize}
      \item \code{int putc(int character, FILE *stream)}
    \end{itemize}

    \pause

    \item Zápis řetězce
    \begin{itemize}
      \item \code{int fputs(const char *str, FILE *stream)}
    \end{itemize}

    \pause

    \item Formátovaný zápis
    \begin{itemize}
      \item \code{int fprintf(FILE *stream, const char *format, ...)}
    \end{itemize}

    \pause

    \item Blokový zápis na binární úrovni
    \begin{itemize}
      \item \code{size_t fwrite (const void* ptr, size_t size, size_t count,
	FILE *stream)}
      \item Zapíše blok bajtů ptr o zadané délce: \code{size * count}
    \end{itemize}
  \end{itemize}

\end{frame}

\begin{frame}[fragile]{Zápis do souboru}{Příklad}

  \begin{block}{Formátovaný zápis po znacích}
  \begin{lstlisting}
#include <err.h>
#include <stdio.h>

struct fruit {
  char sort[32];
  float size;
};

void my_write() {
  struct fruit f = { "Orange", 2.5 };
  FILE *file;

  file = fopen("fruits.txt", "w");
  if (!file)
    err(1, "fopen");
  fprintf(file, "%.2f kg of %s\n", f.size, f.sort);
  fclose(file);
}
  \end{lstlisting}
  \end{block}

\end{frame}

\begin{frame}
  \frametitle{Úkol}

  \heading{Zápis do souboru}
  \begin{enumerate}
    \item Zapište počet řádků v souboru na konec souboru
    \begin{itemize}
      \item I s novým řádkem
    \end{itemize}
    \item Vyzkoušejte
    \item Vypište si soubor
  \end{enumerate}

\end{frame}

\begin{frame}{Změna pozice v souboru}

  \begin{itemize}
    \item Aktuální pozici v souboru lze měnit bez čtení/zápisu
    \item \code{int fseek(FILE *stream, long offset, int whence)}
    \begin{itemize}
      \item \code{offset} -- zadaný offset vzhledem k \code{whence}
      \item \code{SEEK_SET} -- začátek souboru
      \item \code{SEEK_CUR} -- aktuální pozice
      \item \code{SEEK_END} -- konec souboru
    \end{itemize}

    \pause

    \item \code{void rewind(FILE *stream)}
    \begin{itemize}
      \item Přesune aktuální pozici na začátek souboru
    \end{itemize}

    \pause

    \item Zjištění aktuální pozice
    \begin{itemize}
      \item \code{long ftell(FILE *stream)}
    \end{itemize}

  \end{itemize}

\end{frame}

\begin{frame}
  \frametitle{Úkol}

  \heading{Změna aktuální pozice v souboru}
  \begin{enumerate}
    \item Přečtěte a vypište 10 bajtů ze souboru
    \item Přesuňte se na 2. bajt v souboru, přečtěte 10 a vypište
    \item Opakujte 10$\times$ a zvyšujte pozici o 1
    \item Vypište posledních 30 bajtů ze souboru
    \item Vyzkoušejte
  \end{enumerate}

\end{frame}

\begin{frame}{stdin, stdout, stderr}

  \begin{itemize}
    \item Standardní držátka k souboru
    \item Automaticky otevřeny a zavřeny systémem

    \pause

    \item \code{printf() == fprintf(stdout)}
    \item \code{scanf() == fscanf(stdin)}
    \item \code{getchar() == getc(stdin)}
    \item \dots
  \end{itemize}

\end{frame}

